\documentclass[amsmath,amssymb,aps,prb]{revtex4-2}
%\documentclass[a4paper,12pt]{article}
%\usepackage[a4paper, total={6.0in, 8in}]{geometry}
%\usepackage[utf8]{inputenc}

\usepackage{color}
\usepackage{graphicx}
\usepackage{pgfplots,mathtools}
\usepackage{hyperref}
\usepackage{bm}
\usepackage{amsmath}
\usepackage{braket}
\usepackage{slashed}
\usepackage{mathrsfs}
\usepackage[compat=1.0.0]{tikz-feynman}
%\usepackage[super,compress]{cite}


%===============================================
\newcommand{\be}{\begin{equation}}
\newcommand{\ee}{\end{equation}}
\newcommand{\bea}{\begin{eqnarray}}
\newcommand{\eea}{\end{eqnarray}}
\newcommand{\ba}[1]{\begin{array}{#1}}
\newcommand{\ea}{\end{array}}
\newcommand{\nn}{\nonumber}

\newcommand{\ep}{\epsilon}
\newcommand{\om}{\omega}  
\newcommand{\vk}{\vec k}
\newcommand{\vq}{\vec q}
\newcommand{\vp}{\vec p}
%===========================================================
\newcommand{\FB}[1]{\left(#1\right)}
\newcommand{\bp}{\boldsymbol{p}}
\newcommand{\del}{\partial}
%=================================================
%
\begin{document}
%
\title{Shear Viscosity expression for Graphene system in Relaxation time approximation}
%
\author{Cho Win Aung, Thandar Zaw Win, Gaurav Khandal,
Sabyasachi Ghosh}
%
 \affiliation{Indian Institute of Technology Bhilai, GEC Campus, Sejbahar, Raipur 492015, Chhattisgarh, India}
%
\begin{abstract}
We have gone through the detailed microscopic calculation of the shear viscosity of a 2-dimensional graphene system in the 
relaxation time approximation-based kinetic theory framework. After getting its final expressions, we compared
it with the shear viscosity expressions of other possible 2-dimensional as well as 3-dimensional non-relativistic and 
ultra-relativistic fluid systems. The aim of the comparison is to reveal - how their different one-body dispersion relations
affect their many-body fluid properties like shear viscosity and viscosity to entropy density ratio. It is also aimed to
reveal the 3-dimension to the 2-dimension transformation of their mathematical structures. We have numerically
explored the differences in their order of magnitude and dependence on thermodynamical parameters - temperature and chemical
potential. Marking two thermodynamical domains - Dirac fluid and Fermi liquid for a 2-dimensional graphene system, we have noticed
that shear viscosity, entropy density as well as their ratios decrease towards saturated values when one goes from Fermi liquid
to Dirac fluid domain. When one shifts from mili-electron Volt scales of temperature and chemical potential in condensed matter physics
location to their Mega-electron Volt scales in high energy physics location, then the same results may be expected for hot quark matter case,
where the transition from the neutron star to early universe domains may be considered as Fermi liquid to Dirac fluid transition. 



%We have explored the ratio of shear viscosity to energy entropy for the hydrodynamic flow of electrons in a 2D graphene layer. Since the ratio of shear viscosity to entropy density describes the fluid properties and the interactions among the constituents in the hydrodynamic system, we initiated the ideal part of the energy-momentum tensor to get the thermodynamic physical quantities and shear viscosity of electrons. Based on the kinetic theory approach relaxation time approximation, the Boltzmann transport equation for the liquid flow in graphene is used to derive the shear viscosity coefficient. Our calculated result of the ratio is close to the lower bound of the strongly interacting matter. As a result, the hydrodynamic flow of electrons in Dirac fluid can also be explained briefly regarding the constituent carriers in the fluid.
% Interestingly,
% fluidity of quark-gluon plasma remains unaffected by thermodynamical interaction, tuned from LQCD data,
% but not by a magnetic field and dissipation interaction, quantified by collisional relaxation time. In other
% word, one can speculate that the perfect fluid nature of QGP is solely governed by dissipation interaction and 
% probably external magnetic field, produced in heavy ion collision.
\end{abstract}
%
\maketitle
%
\section{Introduction} 
    It is known that the mean free path of charge carriers in metal is generally temperature dependent. The scattering between electron and lattice imperfections (or “disorder”) normally dominates at low temperatures, while electron-phonon scatterings dominate at high temperatures. Concerning these two scattering mechanisms, another possible scattering is electron-electron scattering processes, which are generally less effective in many conventional metals. However, its opposite condition is possible in some specific systems under specific conditions, where one can apply the electron hydrodynamic (eHD) theory. For a long time, condensed matter physicists did not aware of such an opposite phase in materials. Therefore, they used to give less attention to the possibilities of the hydrodynamics behavior of electrons. After the experimental observations of eHD in Refs.~\cite{1Ku:2019lgj,2varnavides2020electron,3PhysRevB.103.125106,4PhysRevB.103.235152,5PhysRevB.103.155128,6PhysRevResearch.3.013290,7DiSante:2019zrd,8PhysRevB.103.115402,9sulpizio2019visualizing,10gallagher2019quantum,11doi:10.1126/science.aau0685,12ella2019simultaneous,13bandurin2018fluidity,18doi:10.1126/science.aad0201,14PhysRevB.98.241304,15jaoui2018departure,16doi:10.1126/science.aac8385,17doi:10.1126/science.aad0343}~, the situation has drastically changed in recent years. See Refs.~\cite{eHD1,eHD2,eHD3} for recent reviews. It is graphene \cite{1Ku:2019lgj,2varnavides2020electron,3PhysRevB.103.125106,4PhysRevB.103.235152,5PhysRevB.103.155128,6PhysRevResearch.3.013290,7DiSante:2019zrd,8PhysRevB.103.115402,9sulpizio2019visualizing,10gallagher2019quantum,11doi:10.1126/science.aau0685,12ella2019simultaneous,13bandurin2018fluidity,18doi:10.1126/science.aad0201}, which is identified as the best known such material, where electron hydrodynamics can be observed. Apart from these recently discovered hydrodynamic properties of electrons in graphene, it was quite famous for its massless nature, concluded from the proportional relation between its energy and momentum. Due to the proportional relation between energy and momentum, electron motion in graphene will not be Galilean-invariant. On the other hand, the relativistic effect of electrons can not also be expected because its velocity ($v_g\approx 10^6$ m/s) is not very close to the speed of light ($c\approx 3\times 10^8$ m/s). Hence, we can’t  claim the Lorentz-invariant property of electron motion. It opens “unconventional” hydrodynamics~\cite{eHD2} as neither non-relativistic hydrodynamics (NRHD) nor relativistic hydrodynamics (RHD) can be applicable. We may call this “unconventional” hydrodynamics as Graphene hydrodynamics (GHD) by imposing that the graphene (G) case has a unique dispersion or energy-momentum relation and is different from the non-relativistic (NR) and relativistic (R) or ultra-relativistic (UR) cases.
%
    Now, whenever fluid dynamics or hydrodynamics comes into the picture, then one dissipation coefficient like the shear viscosity of that fluid, becomes a very important quantity, which is not at all appeared in most of the metals or other condensed matter systems. The present work is aimed at the microscopic calculation of the shear viscosity of this electron fluid in a graphene system, which may be called in short as graphene fluid (GF). When one microscopically calculates the expression of the shear viscosity of GF, it will be different from its standard expression for non-relativistic fluid (NRF) as well as for relativistic fluid (RF) or ultra-relativistic fluid (URF).    
%
    So far, from the best of our knowledge, experimental measurement of this shear viscosity coefficient for GF is missing although experimental community~\cite{1Ku:2019lgj,9sulpizio2019visualizing} observed the Poiseuille's flow pattern of electrons in graphene, which indirectly reflects the existence of the non-zero viscosity. From the theoretical side, we get only Refs.~\cite{24PhysRevLett.103.025301, Vignale}, where microscopic expressions of shear viscosity have been addressed. In this context, one can get a long list of Refs.~\cite{QGP1, QGP2, QGP3, QGP4, QGP5, QGP6, QGP7, QGP8, QGP9, QGP10} (and references therein) for microscopical estimations of shear viscosity for relativistic quark and hadronic matter, expected in high energy heavy ion collision experiments. Grossly two classes of frameworks - (1) Kinetic theory approach with relaxation time approximation (RTA)~\cite{QGP1, QGP2, QGP3, QGP4, QGP5, QGP6, QGP7} and Kubo framework~\cite{QGP8, QGP9, QGP10} are adopted by the heavy ion physics community. Both frameworks have similar structure at the final level expressions for shear viscosity coefficients with two main components. One carries interaction information, called relaxation time, and the remaining part may be called as thermodynamic phase-space of shear viscosity coefficient, which will be function of temperature and chemical potential. If we analyze the shear viscosity expression of graphene also from Ref.~\cite{24PhysRevLett.103.025301}, then we can identify these two components. Present work has zoomed in this structure via a systematic calculation of this shear viscosity of GF in RTA methods and compared with corresponding structures for NRF and URF. Here, one of our aims is to compare the thermodynamic phase-space component of shear viscosity coefficient for these three cases - G, NR, and UR.  
%
    After knowing the lower bound conjecture of shear viscosity to entropy density $(\eta/s)$ as $\hbar/(4\pi k_B)$ or $1/(4\pi)$ (in natural unit)~\cite{25Kovtun:2004de}, scientific communities are curious to know those strongly coupled systems, which are close to that bounds. Experimentally, the RF, like quark and hadronic matter, produced in high energy heavy ion collision experiments and NRF like cold atom systems~\cite{27Sch_fer_2009} are identified as those strongly coupled systems. On the other hand, GF may also belong to that category according to the theoretical prediction from Ref.~\cite{24PhysRevLett.103.025301}, which is considered as reference point for tuning our results. So, the present article will not intend to add any new content on strongly coupled properties, rather its main goal will intend to find the differences among GF, NRF, URF in terms of expressions and estimations of shear viscosity. We believe that it was missing in the literature and very important to address.
    %The ratio of shear viscosity to entropy density $(\eta/s)$ can be used to characterize a given fluid as perfect. From the string theory calculation, the ratio is equivalent to a universal constant value of $\hbar/(4\pi k_B)$ \cite{25Kovtun:2004de}, which can be regarded as the lower bound of a wide range of systems or KSS bound. It is a more meaningful quantity than the viscosity; it directly measures the fluid property \cite{26Dey:2019axu}. Therefore, our attention comes to the Graphene system to explore the different dispersion relations or energy-momentum relation of electron flow on Graphene, where we calculated the shear viscosity to entropy density ratio, comparing it with the non-relativistic system. Since graphene is composed of a single layer of carbon atoms arranged in a two-dimensional (2D) honeycomb lattice, we have paid attention to the two-dimensional description of shear viscosity and its related thermo-dynamical properties. 
%
    The article is organized as follows. In the next section Sec.~(\ref{sec: Form}), the RTA calculation of shear viscosity $\eta$ and entropy density $s$ calculations of GF for 2D case is addressed in detail by mentioning the other cases like 3D-GF, 3D-NR, 3D-UR, 2D-NR, 2D-UR. In Sec.~(\ref{sec: Result}), the comparative results of $\eta$, $s$ and $\eta/s$ of different cases are discussed. At the end, our findings are summarised in Sec.~(\ref{sec: Sum}) with some conclusive bullet points.
    %By using that “unconventional” hydrodynamical description, the present work has attempted to describe a dissipation picture of electron hydrodynamics in the presence of a magnetic field. We mainly concentrate on calculating different shear viscosity components in the presence of an external magnetic field.
%
    \section{Formalism}
    \label{sec: Form}
    Let's start our formalism from energy-momentum tensor ($T^{\mu\nu}$), as practised for RF like quark and hadronic matter. Here, we will go for GF calculation, so reader should have to be careful on some particular steps, where it is different from the RF case. Showing these differences is one of the core agenda of the present article. Although, reader can find similarities between most of steps of GHD of GF and RHD of RF. The $T^{\mu\nu}$ has two parts - ideal part $T_0^{\mu\nu}$, related to the knowledge of thermodynamics and dissipative part $T_D^{\mu\nu}$, related to the different dissipation processes. So,
    \bea
    T^{\mu\nu} &=& T_0^{\mu\nu} + T_D^{\mu\nu}~. 
    \eea
    In this dynamic picture of fluid, ideal energy-momentum tensor and electron number flow can be expressed in macroscopic form as,
    \bea
    T_0^{\mu\nu} &=&  \ep \frac{u^\mu u^\nu}{v_g^2} - P\left(g^{\mu\nu} -\frac{u^\mu u^\nu}{v_g^2}\right)~,
    \nn\\
    N_0^\mu &=& n \frac{u^\mu}{v_g}~,
    \label{eq2}
    \eea
    in terms of the building blocks - energy density $\epsilon$, pressure $P$, number density $n$, fluid (element) velocity $u^\mu$ and metric tensor $g^{\mu \nu}$. Here, four-velocity $u^{\mu}=\gamma_g (v_g, \vec{u})$ for GHD is designed by following the four velocity structure $u^{\mu}=\gamma (c, \vec{u})$ for RHD as done in Refs.~\cite{eHD2}. One can notice that speed of light $c$ in RHD is replaced by graphene Fermi velocity $v_g$ in GHD. So, Lorentz factor $\gamma= 1/\sqrt{1-u^2/c^2}$ in RHD is also converted into $\gamma_g= 1/\sqrt{1-u^2/v_g^2}$ in GHD. In static limit (${\vec u}\rightarrow 0$), four velocity, $u^{\mu}=\gamma_g (v_g, \vec{u}) \rightarrow u^{\mu}=\gamma_g (v_g, 0)$ and $\gamma_g= 1/\sqrt{1-u^2/v_g^2} \rightarrow 1$. So, Eq.~(\ref{eq2}) provides a static electron number flow $N^\mu_0\equiv n$, and static energy-momentum tensor, 
    \begin{equation}
    T_0^{\mu\nu} \equiv  \begin{pmatrix}
    \epsilon &  0   &  0 & 0\\
     0  &   P &  0 & 0\\
     0  &  0  &  P &  0\\
     0  &  0  &  0 &  P\\
     \end{pmatrix}~,
%
    \end{equation}
    which reflects  the standard static fluid aspect like Pascal's law. The macroscopic quantities $T_0^{\mu \nu}$ and $N^\mu_0$ can be expressed in terms of the microscopic quantities; four-momentum ($p^{\mu}$) and four-velocity ($v^{\mu}$) of electrons as,
    \begin{equation}
    T_0^{\mu\nu} = N_s\int \frac{d^3\vec{p}}{(2\pi)^3}p^\mu v^\nu f_0~,
    \label{rohan}
    \end{equation}
    and
    \begin{equation}
    N_0^{\mu} = N_s\int \frac{d^3\vec{p}}{(2\pi)^3} v^\mu f_0~,
    \label{N_mic}
    \end{equation}
    where $N_s=2$ is spin degeneracy factor of electron and $f_0$ is its Fermi-Dirac (FD) distribution function $f_0=1/\{\exp{(\beta(E-\mu))}+1\}$. Here, $\beta=1/(k_B T)$ and $\mu$ are the thermodynamic parameter and the chemical potential of the system, respectively. From these microscopic expressions of the energy-momentum tensor and electron current, given in Eqs.~(\ref{rohan}) and (\ref{N_mic}), we can write the energy density $\ep$, pressure $P$ and number density $n$ for 2D graphene (G) case, which is addressed briefly in next subsection. We follow natural unit $\hbar=c=k_B=1$ during the calculation.
%
    \subsection{Entropy density in two-dimensional Graphene}

    For Graphene, the dispersion relation is given by
    \begin{equation}
    E = pv_g~.
    \label{nanu}
    \end{equation}
    The total number of fermions at any value of temperature is given by,
    \begin{equation}
    N = \int_0^{\infty}D\left(E\right)dE f_0~,
    \label{3dnd}
    \end{equation}
    where $D\left(E\right)dE$ is number of energy states in energy range $E$ to $E+dE$.
%
    After plugging the value of $D\left(E\right)dE$ (see Appendix~\ref{App_A}) in the above Eq.~(\ref{3dnd}) and $f_0$, the total number of electrons in graphene is,
    \begin{align*}
    &N = N_s \frac{2\pi a}{\left(2\pi\right)^2 v_g^2} \int_{0}^{\infty}  \frac{E}{A^{-1}e^{\beta E} + 1} dE~,
    \end{align*}
    %where $N_s=2$ is the spin degeneracy of the electron.
    %that arises due to the spin up and spin down energy state of the same energy and which is $N_s=2s+1$ where $s$ is the spin of a particle. In the case of graphene, the Brillouin zone has two in-equivalent Dirc points ("valleys"), each with two degrees of freedom. So in graphene, another two degeneracy factor exists, known as valley index, but we are not adding this in our calculation as our interest to compare G case with NR and UR cases. 
    where $A=\exp{(\beta \mu)}$ is the fugacity of the system. After converting this integral into the Fermi integral function (see Appendix~\ref{App_B}), we get the expression of number density as 
    \begin{equation}
    n_g^{2D} = \frac{N}{a} = \frac{N_s}{2 \pi v_g^2} f_2\left(A\right) T^2~.
    \end{equation}
%
    Now from the Eq.~(\ref{rohan}), the energy density for (2D) graphene system is,

    \begin{equation}
    \epsilon_g^{2D} = T_0^{0 0} =  N_s \int \frac{d^2 p}{\left(2\pi\right)^2}\left(E \right) f_0.
    \end{equation}
%
    After using the graphene dispersion relation and plugging the value of $f_0$, we get
    \begin{equation}
    \epsilon_g^{2D} = \frac{1}{\pi v_g^2} \int_{0}^{\infty}  \frac{E^2}{A^{-1}e^{\beta E} + 1} dE
    \end{equation}
    and after replacing this integral in terms of Fermi integral function, the final expression of the energy density is given by
    \begin{equation}
    \epsilon_g^{2D} = \frac{N_s}{\pi v_g^2} f_3\left(A\right) T^3~.
    \end{equation}
    Now again from the Eq.~(\ref{rohan}), the pressure can be expressed as 
    \begin{equation}
    P_g^{2D} = T_0^{11} = N_s \int \frac{d^2 p}{\left(2\pi\right)^2}\left(\frac{E}{2} \right) f_0~,
    \end{equation}
    since $\vec{p}_x \vec{v}_x \approx \frac{|\vec{p}|}{\sqrt{2}} \frac{|\vec{v}_g|}{\sqrt{2}}=\frac{E}{2}$.
    After solving this expression as similar to the energy density, we get
    \begin{equation}
    P_g^{2D} = \frac{N_s}{2 \pi v_g^2} f_3\left(A\right) T^3.
    \end{equation}
%
    In terms of number density, energy density, and pressure, we can write the entropy density from the Euler thermodynamic relation:
    \begin{equation}
    s = \frac{S}{V} = \frac{\epsilon + P - \mu n}{T}~.
    \label{viral}
    \end{equation}
%
    After substituting the value of energy density ($\epsilon_g^{2D}$), pressure ($P_g^{2D}$), and number density ($ n_g^{2D}$) in Eq.~(\ref{viral}), we get
    \begin{equation}
    s_g^{2D} = \frac{N_s}{2 \pi v_g^2} T^2 \Big[3f_3\left(A\right) - \frac{\mu}{T} f_2\left(A\right) \Big]~.
    \label{pure}
    \end{equation}
%
 %%%%%%%%%%%%%%%%%%%%%%%%%%%%%%%%%%%%%%%%%%%%%%%%%%%%%%%%%%%%%%%%%%%%%%%%%
    % \newpage
    %\begin{figure}[!h] 
    %  	\centering 
    %    	\includegraphics[scale=0.25]{Th_eta_mu.eps}
    %    	\includegraphics[scale=0.25]{Th_eta_T.eps}
    %    	\caption{The shear viscosity at $T=0$ and chemical energy $\mu$ (Left) and the shear viscosity at $\mu=0$ and chemical energy (Right) of electron hydrodynamics flow in different categories of physics.} 
    %    		\label{fig: ---}
    %\end{figure} 
%
 %%%%%%%%%%%%%%%%%%%%%%%%%%%%%%%%%%%%%%%%%%%%%%%%%%%%%%%%%%%%%%%%%%%%%%%%
    \subsection{Shear Viscosity in two-dimensional Graphene}
%\hfill \break
%%%%%%%%%%%%%%%%%%%%%%%%%%%%%%%%%%%%%%%%%%%%%%%%%%%%%%%%%%%%%
    Next, let us come to dissipative part of $T_D^{\mu\nu}$, where only shear stress tensor $\pi^{\mu\nu}$ will be considered for calculating shear viscosity coefficient ($\eta$). Macroscopic relation $\pi_{\mu\nu}=\eta \mathcal{U}_{\mu\nu}$ defines that the shear viscosity $\eta$ is a basically proportional constant between shear stress $\pi^{\mu\nu}$ and velocity gradient $\mathcal{U}_{\mu\nu}=\frac{1}{2} (\partial_{\mu} u_{\nu}+\partial_{\nu} u_{\mu})$.
    Usually, greek index like $\mu\equiv(0,i)$ takes values 0 for the temporal component and $i=1,2,3$ for the spatial component for 3D system but here for 2D system, we will consider $\mu\equiv(0,i=1,2)$ because z-component $i=3$ will not be considered.
%
    The microscopic theory describes shear stress tensor in terms of particle velocity $v$ and momentum $p$ as,
%
    \be
    \pi_{\mu\nu}=N_s \int \frac{d^2 \vec{p}}{(2 \pi)^2} p_{\mu} v_{\nu} \delta f_\eta~,
    \label{pij_micro}
    \ee
%
    where we are assuming that equilibrium distribution function $f_0$ gets a small deviation $\delta f$, which can be considered as first-order Taylor series expansion equilibrium distribution function:
%
    \bea 
    \delta f &\propto& \frac{\partial f_0}{\partial E}
    =\phi_\eta \frac{\partial f_0}{\partial E}~,
    \nn\\
    &=&\mathcal{A}^{\mu\nu} \mathcal{U}_{\mu\nu} \frac{\partial f_0}{\partial E}~.
    \eea
%
    Considering the relation $v_{\nu}=(E/p^2) p_{\nu}$, macroscopic $\pi_{\mu\nu}=\eta \mathcal{U}_{\mu\nu}$ and microscopic Eq.~(\ref{pij_micro}) can be connected as, 
%
    \be
    \pi_{\mu\nu}=\eta ~\mathcal{U}_{\mu\nu}= N_s \int \frac{d^2 \vec{p}}{(2 \pi)^2} \left(\frac{E}{p^2}\right)p_{\mu} p_{\nu} \mathcal{A}^{\alpha\beta} \mathcal{U}_{\alpha\beta} \frac{\partial f_0}{\partial E}.
    \ee
%
    The four-momentum of an electron can be defined as $p^\mu=(E/v_g, \vec{p})$ in unconventional notation.
%
    Considering energy as a static limit of $p^\nu u_\nu$, we can write FD as,
%
    \be
    f_0=\frac{1}{exp{\left( \frac{p^\nu u_\nu \mp \mu(x)}{T(x)}\right)+1}}~.
    \label{a1}
    \ee
    Here, we have to consider the local thermalization concept, where thermodynamical quantities $T(x)$, $\mu(x)$ as well as fluid velocity $u^{\mu}(x)$ are assumed to be functions of $x \equiv x^{\mu}=(x^0, x^i)$.
    %For eHD case,  $u^\mu=\gamma_g(v_g, \vec{u})$ will be transformed to $u^\mu=(v_g, \vec{0})$ in fluid rest frame, so $u^\mu p_\mu$ will be converted to $\omega$.
%
    To find the unknown coefficient $\mathcal{A}^{\alpha\beta}$, we will use Boltzmann transport equation (BTE)
    \bea
    \frac{\partial f}{\partial t}+v^\mu \frac{\partial f}{\partial x^\mu}+F^\mu \frac{\partial f}{\partial p^\mu}=\left( \frac{\partial f}{\partial t}\right)_{Col}~,
    \eea
    where $\Big( \frac{\partial f}{\partial t}\Big)_{Col}$ is the collision term that leads the system out of equilibrium. $F^\mu $ is represented as all external forces, and $v^\mu$ is the velocity of the fluid particles.
%
    Using velocity expression in terms of $E$ and $p$ for graphene, $v^\mu=(\frac{E}{p^2})p^\mu$, we get:
%
    \bea
    \left(\frac{E}{p^2}\right)p^\mu \partial_\mu f=\left( \frac{\partial f}{\partial t}\right)_{Col}~,
    \eea
    where we ignore $\frac{\partial f}{\partial t}$ and $F^\mu \frac{\partial f}{\partial p^\mu}$ as they will not contribute in shear dissipation. Using the Relaxation Time Approximation (RTA) method, the collision term can be considered as,
%
    \be
    \left( \frac{\partial f}{\partial t}\right)_{Col}=-\frac{\delta f}{\tau_c}~,
    \ee
%
    where $\tau_c$ is the relaxation time. Putting $f\approx f_0$ in the left-hand side (lhs) of BTE,
%
    \bea
    \left(\frac{E}{p^2}\right) p^\mu  \partial_\mu f_0=-\frac{\delta f}{\tau_c}~.
    \label{fo}
    \eea
%
    Using the Eq.~(\ref{a1}), the lhs of the above Eq.~(\ref{fo}) can be expanded as, 
%
    \bea
    \left( \frac{E}{p^2}\right) p^\mu \partial_\mu f_0 &=& -f_0(1-f_0) \Big[\left(\frac{E}{p^2}\right)\frac{p^\mu p^\nu}{T} \partial_\mu u_\nu (x) \Big]~,
    \nn\\
    &=&-f_0(1-f_0)  \left(\frac{E}{p^2}\right)\frac{p^\mu p^\nu}{2T} ( \partial_\mu u_\nu +\partial_\nu u_\mu)~,
    \nn\\
    &=&-f_0(1-f_0) \left(\frac{E}{p^2} \right)\frac{p^\mu p^\nu}{T} \mathcal{U}_{\mu \nu}~,
    \eea
%
    %Using the identity, $\frac{\partial_\mu T}{T}=u^\alpha \partial_\alpha u_\mu-c^2_s u_\mu \partial_\alpha u^\alpha$, the above equation becomes as,
%
    and right hand side (rhs) of Eq.~(\ref{fo}) can be written as
    \begin{eqnarray}
    -\frac{\delta f}{\tau_c}&=& \frac{ f_0(1-f_0)}{T} \frac{1}{\tau_c} \mathcal{A}^{\mu \nu} \mathcal{U}_{\mu \nu}
    \end{eqnarray}
    %%-\frac{2}{3} \Delta_{\mu \nu} \partial_\alpha u^\alpha \Big)\\
    %&-& (\frac{\omega}{p^2})\frac{p^\mu p^\nu}{3T} \Big( \Delta_{\mu \nu}+ 3c^2_s u_\mu u_\nu   \Big) \mp (\frac{\omega}{p^2}) p^\mu \partial_\mu \Big( \frac{\mu(x)}{T(x)} \Big) \Big]\\
    %Since the force term vanishes, the Eq.(\ref{fo}) becomes,
%
%\begin{eqnarray}
    %    -\Big( \frac{p^\mu u_\mu}{E}\Big) \frac{\delta f}{\tau_c} &=&f_0(1-f_0)\Big( \frac{\vec{p}.\vec{u}}{E T}\Big) \Big( -\frac{E}{p^2} p_\mu p_\nu \tau_c\Big)\frac{1}{\tau_c} \mathcal{U}^{\mu \nu}~,\\
    %    &=& \frac{\vec{p}.\vec{u}}{E T} \frac{1}{\tau_c} \mathcal{A}_{\mu \nu} \mathcal{U}^{\mu \nu} f_0(1-f_0)
    %\end{eqnarray}
%
    So, equating lhs and rhs of Eq.~(\ref{fo}), we get $\mathcal{A}^{\mu \nu}=(-\frac{E}{p^2} p^{\mu} p^{\nu} \tau_c)$. Transforming temporal+spatial to only spatial index, we can write the shear stress tensor as
%
    \begin{eqnarray}
    \pi_{ij}&=& \eta \mathcal{U}_{ij}~,\\
    &=& N_s \int \frac{d^2 \vec{p}}{(2 \pi)^2} \left( \frac{E}{p^2}\right)^2 \tau_c (p_i p_j p_k p_l) \mathcal{U}^{kl} \beta f_0(1-f_0)~,\\
    &=& \frac{N_s}{8} \int \frac{d^2 \vec{p}}{(2 \pi)^2} E^2 \tau_c \beta f_0(1-f_0) \mathcal{U}_{ij}~,
    \end{eqnarray}
%
    where, we used $<p_i p_j p_k p_l>=\frac{\vec{p}^4}{8} \big(\delta_{ij} \delta_{kl}+\delta_{ik}\delta_{jl}+\delta_{il} \delta_{jk} \big)$ (see Appendix~\ref{App_D}), and we have the two Eqs;
 %Now from the equation (\ref{nanu1}), the above expression can be written as,
 %\begin{equation}
 %    \eta= \frac{N_s}{8} \int \frac{d^2 \vec{p}}{(2 \pi)^2} E^2 \tau_c  \beta f_0(1-f_0)~.
 %    \label{GEOSV3D}
 %\end{equation}
    \begin{equation}
    \beta f_0 \left(1-f_0\right) = -\frac{\partial f_0}{\partial E}~,
    \label{waw1}
    \end{equation}
    and
    \begin{equation}
    \left(-\frac{\partial f_0}{\partial E}\right) = \beta \frac{ e^{\beta \left(E - \mu\right)}}{\left(e^{\beta \left(E - \mu\right)} + 1\right)^2} =  \frac{\partial}{\partial \mu} \left( \frac{1}{e^{\beta \left(E - \mu \right)} + 1}\right)~.
    \label{waw2}
    \end{equation}
 Finally, the expression for shear viscosity is,
%
    \begin{equation}
    \eta= \frac{N_s}{8} \int \frac{d^2 \vec{p}}{(2 \pi)^2} E^2 \tau_c \beta f_0(1-f_0)~.
    \label{GEOSV3D}
    \end{equation}
%
    After using the Eqs.~(\ref{waw1}) and (\ref{waw2}) and converting the momentum terms into energy using dispersion relation (\ref{nanu}), the Eq.~(\ref{GEOSV3D}) becomes
    \begin{equation}
    \eta = \frac{N_s }{16 \pi v_g^2} \tau_c \frac{\partial}{\partial \mu } \int_0^{\infty} \frac{E^3}{A^{-1} e^{\beta E} +1} dE~.
    \end{equation}
%
    After solving this integration by using the identity of the  Fermi integral function, the expression of shear viscosity for 2D graphene (using subscript and superscript notation to distinguish the expressions of different systems) is
    \begin{equation}
    \eta_g^{2D} = \frac{3 N_s }{8 \pi v_g^2}\tau_c f_3\left(A\right) T^3~.
    \label{gec1}
    \end{equation}
%
    Now, on taking the ratio of the shear viscosity (\ref{gec1}) and entropy density (\ref{pure}), we get 
    \begin{equation}
    \frac{\eta_{g}^{2D}}{ s_{g}^{2D}} = \frac{3 }{4}\tau_c \Bigg[3f_3\left(A\right) - \frac{\mu}{T} f_2\left(A\right) \Bigg]^{-1} f_3\left(A\right) T~.
    \label{2dg}
    \end{equation}
%
    After doing a similar way of calculation, the expressions of entropy density, shear viscosity, and the ratio of shear viscosity and entropy density for a non-relativistic electron fluid (i.e. $E=p^2/(2m)$) 2-dimensional system are given by
    \begin{align}
    & s_{NR}^{2D} = \frac{N_s m T}{2 \pi } \Bigg[2f_2\left(A\right) - \frac{\mu}{T} f_1\left(A\right) \Bigg]~,
    \label{pure1}\\
    & \eta_{NR}^{2D} = \frac{N_s m }{8 \pi}\tau_c f_2\left(A\right) T^2~,
    \label{e_NR}\\
    & \frac{\eta_{NR}^{2D}}{ s_{NR}^{2D}} = \frac{1}{4} \tau_c \Bigg[2f_2\left(A\right) - \frac{\mu}{T} f_1\left(A\right) \Bigg]^{-1} f_2\left(A\right) T~.
    \label{es_NR}
    \end{align}
    Most of the fluids or liquids (e.g., water) used in our daily life follow non-relativistic fluid dynamics, whose fluid constituent particles obey $E=p^2/(2m)$ dispersion relation. However, for the purpose of comparing, we may assume a hypothetical 2D NR system showing fluid behavior, which may be difficult to be found in the real world. By this comparison (given details in the results section), our aim is to encourage the scientific community to use the expressions of G-case, given in Eqs.~(\ref{gec1}), (\ref{2dg}) instead NR-case, given in Eqs.~(\ref{e_NR}), (\ref{es_NR}) when they are describing eHD in graphene system.
%

    If we consider graphene as a 3-dimensional (3D) system, following the dispersion relation $E=pv_g$, then it may be again a hypothetical example but a good example for comparison purpose. Modifying our above calculation with replacement of $\int d^2p\rightarrow \int d^3p$ and $p_xv_x\approx \frac{pv_g}{3}=\frac{E}{3}$, we get the expressions of entropy density, shear viscosity, and their ratio as
    \begin{align}
    & s_g^{3D} = \frac{N_s T^3}{\pi^2 v_g^3} \Bigg[4f_4\left(A\right) - \frac{\mu}{T} f_3\left(A\right) \Bigg]~,
    \label{sg_3D}
    \\
    & \eta_{g}^{3D} = \frac{4 N_s }{5 \pi^2 v_g^3} \tau_c f_4\left(A\right) T^4~,
    \label{eg_3D}
    \\
    & \frac{\eta_g^{3D}}{ s_g^{3D}} = \frac{4}{5} \tau_c \Bigg[4f_4\left(A\right) - \frac{\mu}{T} f_3\left(A\right) \Bigg]^{-1}  f_4\left(A\right) T~.
    \label{esg_3D}
    \end{align}
    %Where $\sigma$ is the cross-section of 3-dimensional graphene, the area perpendicular to the velocity of electrons.\\
%
    Now we have a 3-dimensional non-relativistic (3D NR) system of fermions. After applying the same methodology to this system, we get all the expressions of entropy density, shear viscosity, and the ratio of shear viscosity to entropy density,
%
    \begin{align}
    & s_{NR}^{3D} = N_s \left( \frac{m}{2 \pi }\right)^{\frac{3}{2}} T^{\frac{3}{2}} \Bigg[\frac{5}{2} f_{\frac{5}{2}}\left(A\right) - \frac{\mu}{T} f_{\frac{3}{2}}\left(A\right) \Bigg]~,
    \label{sNR_3D}
    \\
    & \eta_{NR}^{3D} = \frac{N_s}{4} \left( \frac{m}{2 \pi }\right)^{\frac{3}{2}} \tau_c f_{\frac{5}{2}}\left(A\right) T^{\frac{5}{2}}~,
    \label{eNR_3D}
    \\
     & \frac{\eta_{NR}^{3D}}{ s_{NR}^{3D}} = \frac{1}{4} \tau_c \Bigg[\frac{5}{2} f_{\frac{5}{2}}\left(A\right) - \frac{\mu}{T} f_{\frac{3}{2}}\left(A\right) \Bigg]^{-1} f_{\frac{5}{2}}\left(A\right) T~.
     \label{esNR_3D}
    \end{align}
    This 3D NR system, showing fluid behavior, can be applicable for most of the fluids or liquids (e.g. water) used in our daily life.  
    One can consider the above shear viscosity, entropy density, and their ratio for the water molecule, where water molecule obeying NR dispersion relation, $E=p^2/(2m)$ with effective mass $m$, but that will be not a good example to compare with the same expressions for eHD in graphene case. So, we can again consider a hypothetical example - 3D eHD NR case.

    We can also compare the above expressions for 2D, 3D eHD of G and NR cases with the same for ultra-relativistic (UR) case. A good example is hot QGP, where RHD can be applicable. According to the latest understanding~\cite{VR_Rev}, RHD is quite successful in describing QGP phenomenology. Again to make our comparison on equal footing, we will consider the hypothetical case of 2D, 3D UR electron fluid.
    If the Fermi velocity of electrons in graphene $v_g$ is replaced by the factor $1$ (as $c=1$ in natural unit), then all expressions of the G-case will be converted to corresponding expressions of UR case.
%
%%%%%%%%%%%%%%%%%%%%%%%%%%%%%%%%%%%%%%%%%%%%%%%%%%%%%%%%%%%
%\newpage
    \section{Results}
    \label{sec: Result}
%
    After addressing the final expressions of $\eta$, $s$, and its ratio for different systems like 2D, 3D NRF, GF and URF in the formalism section, here we will discuss their numerical estimations through different graphs. 
%
    \begin{figure}
    \centering
    \includegraphics[scale=0.47]{mu_T_plot.eps}
    \caption{Location of condensed matter physics (CMP) domain and the high energy physics (HEP) domain in $T$-$\mu$ diagram.}
    \label{fig:mu_T}
    \end{figure}



    Let us first come to the entropy density results. In the early universe scenario, a hot quark-gluon plasma (QGP) state around temperature $T=400$ MeV or $T=700$ MeV and zero quark chemical potential ($\mu=0$) is expected just after a few micro-second from the big bang. Due to the very high temperature of the medium, the constituent particle average momenta become so large that we can ignore its mass term, and it can be considered a UR case. UR case is famous for photon gas or black body radiation example where internal energy density or Intensity (they have connecting relation) follows $T^4$ law, popularly known as Stefan-Boltzmann (SB) law. QGP thermodynamics at high temperatures reaches that SB limits. Eq.~(\ref{sg_3D}) can be converted to UR case by replacing $v_g=c=1$ and by putting $\mu=0$, we can get SB limit expression for 3D case,
    \be 
    s_{SB} = \frac{N_s T^3}{\pi^2 } \Bigg[4\times \Big\{\frac{7}{8}\zeta_4\Big\} \Bigg]~,
    \label{sSB_3D}
    \ee 
    where Fermi integral function can be converted into Riemann zeta function $f_4=\frac{7}{8}\zeta_4$ for $\mu=0$ condition, following the general relation $f_n=\Big(1-\frac{1}{2^{n-1}}\Big)\zeta_n$. For the QGP case, quark degeneracy factors will have to be put into the $N_s$, and gluon contribution must be added separately. Since two-flavor quark has a degeneracy factor of 24 and gluon has a degeneracy factor of 16, so massless QGP entropy density or SB limits of QGP will be 
    \be 
    s^{QGP}_{SB} = \frac{24 T^3}{\pi^2 } \Big[4\times \Big\{\frac{7}{8}\zeta_4\Big\} \Big]+\frac{16 T^3}{\pi^2 } \Big[4\times \zeta_4 \Big].
    \label{s_QGP}
    \ee 
    When we plan to compare graphene entropy density with this SB limit, we have to understand that the temperature range (few hundred mega electron Volt (MeV), which is equivalent to $10^{12}~^0$K ) of QGP is too much larger than temperature range (1-23 milli electron Volt (meV), which is equivalent to 15-300 $^0$K ) of graphene system. Fig.~(\ref{fig:mu_T}) has addressed nicely about these two domains. It is basically $T$ vs. $\mu$ plots in log scale to cover a broad band of $T$ and $\mu$ range. We marked the condensed matter physics (CMP) domain, covering  $T\approx 1-23$ meV and $\mu\approx 0-10$ eV. We know that metal Fermi energy remains within the range $\mu=2-10$ eV, which is marked as yellow. Unlike metal, graphene system Fermi energy can be changed via doping methods, and its $\mu/T\ll 1$ and $\mu/T\gg 1$ domains are called Dirac fluid (DF) or Dirac liquid (DL) and Fermi liquid (FL) domains, respectively, marked by arrows in Fig.~(\ref{fig:mu_T}). Similar to DF and FL domains for electrons, we may call early universe QGP as DF domain of quark and quark matter, expected in the core of neutron star as FL domain of quark. A rectangular domain within $T=1-400$ MeV and $\mu=0-1000$ MeV is marked as high energy physics (HEP) domain for quark. Reader can easily visualize the gap between CMP and HEP domains.

%
    \begin{figure}
    \centering
    \includegraphics[scale=0.25]{Fig3_Gen_s_ssb_muT_3D.eps}
    \includegraphics[scale=0.25]{Fig3_Gen_s_ssb_muT_2D.eps}
    \caption{The ratio of entropy density in different domains to $s_{SB}$ with $\mu/T$ (a) in the 3D case and (b) in the 2D case respectively}
    \label{fig:s}
    \end{figure}
%
    After realizing the scale gap in $T$-$\mu$ plane between URF of quark and GF of electrons, one should understand that we must consider a hypothetical electron URF to make an equal footing comparison. Within the temperature ($T=0-0.023$ eV) and chemical potential ($\mu=0-10$ eV) range, entropy density of URF,
    \be 
    s_{UR}^{3D} = \frac{N_s T^3}{\pi^2} \Bigg[4f_4\left(A\right) - \frac{\mu}{T} f_3\left(A\right) \Bigg]~,
    \ee 
    has to be plotted with a normalization by $s_{SB}$, given in Eq.~(\ref{sSB_3D}). This normalized estimation is sketched by the blue dotted line in the left panel of Fig.~(\ref{fig:s}), which shows that $s_{UR}^{3D}\approx s_{SB}$ in the domain $\mu/T\ll 1$, as expected. Interestingly, we noticed that the main $\mu/T$ dependence in entropy density is coming beyond the $\mu/T=1$. Reader can understand that the terms with Fermi integral function are the main source of $\mu/T$ dependence. Next, using Eq.~(\ref{sg_3D}), the graphene entropy density for the Fermi velocity $v_g=0.006$ is plotted (red solid). From Ref.~\cite{vg_Nature}, we can get knowledge about a broad range of Fermi velocity $v_g=1-3\times 10^6$ m/s or $v_g=0.003-0.01$ (in natural unit) in graphene system. As charge career density or $\mu$ decreases, $v_g$ will increase and approach towards Dirac fluid (DF) or strongly coupled electron-electron domain. We have considered in-between constant values $v_g=0.006$. We can understand that the $\mu/T$ dependence of entropy density for URF and GF are the same but $GF\gg URF$ due to the $1/v_g^3\approx 5\times 10^6$ term. Next, we use Eq.~(\ref{sNR_3D}) to draw the entropy density of NRF to plot (green dash line) in the left panel of Fig.~(\ref{fig:s}). Reader can understand its different trend of $\mu/T$ dependence for NRF is because of the term $\Big[\frac{5}{2} f_{\frac{5}{2}}\left(A\right) - \frac{\mu}{T} f_{\frac{3}{2}}\left(A\right) \Big]$.

    A similar trend we can notice for 2D case with similar ranking URF $\ll$ GF $\ll$ NRF. Only for the transition from 3D to 2D, their orders of magnitude are shifted toward lower values. 

    Next, let us come to the shear viscosity results. Here also, we can expect SB limit type simple expression for UR case at $\mu=0$:
    \be 
    \eta_{SB} = \frac{4 N_s }{5 \pi^2} \tau_c \frac{7}{8}\zeta_4 T^4~,
    \label{eSB_3D}
    \ee 
    from the general $\eta(T,\mu)$ expression for URF:
    \be 
    \eta_{UR}^{3D} = \frac{4 N_s }{5 \pi^2} \tau_c f_4\left(A\right) T^4~,
   \label{eUR_3D}
    \ee 
    by putting $v_g=c=1$ in Eq.~(\ref{eg_3D}).
    For massless QGP at $\mu=0$ case, by replacing degeneracy factors of quark and gluons in $N_s$, we get 
    \be 
    \eta^{QGP}_{SB} = 24\Big[\frac{4}{5\pi^2 } \tau_c\times \Big\{\frac{7}{8}\zeta_4\Big\}T^4 \Big]+16\Big[\frac{ 4}{5\pi^2 } \tau_c\times \zeta_4 T^4 \Big]~.
    \ee
    Again, this QGP is a realistic example of URF but for comparison, we have to consider electron URF. When we plan to compare the shear viscosity of URF, GF, and NRF, then we should use Eqs.~(\ref{eUR_3D}), (\ref{eg_3D}), (\ref{eNR_3D}) and for SB limit, we will use Eq.~(\ref{eSB_3D}). Similar to normalized entropy density by its SB limit in Fig.~(\ref{fig:s}), we have plotted normalized shear viscosity by its SB limit in Fig.~(\ref{fig:e}), where 3D and 2D estimations are plotted in left and right panels respectively. Shear viscosity expression carries two kinds of information. One is relaxation time $\tau_c$, and another is the remaining thermodynamic phase-space part as a function of $T$ and $\mu$. During the normalization, $\tau_c$ information is canceled, and we can only see their thermodynamic phase-space part of shear viscosity. Interestingly, it follows a similar trend to other thermodynamical quantities like entropy density - which shows two types of $\mu/T$ dependence in the domain $\mu/T\ll 1$ and $\mu/T\gg 1$, which are commonly assigned as DF and FL.


    %To trace the suspect of why the ratios of the shear viscosity to entropy density have some changes in the F.L region, we also analyze the shear viscosity and the energy entropy density due to electron flow on Graphene separately, comparing these ratios with the respective values of $\eta$ and $s$ in different cases in the Fig.~(\ref{fig:my_label3}) and in the Fig.~(\ref{fig:my_label4}). Consequently, a classical ultra-relativistic gas with zero chemical potential can be regarded as Wein and Planck gas from the Boltzmann distribution function; $\eta_{UR}$ is equivalent to $\eta_{SB}$, i.e., they both obey Stefan-Boltzmann law \cite{28sym11020249}. The ratios of the shear viscosity have the same trend as we see $\eta/s$ in Fig.~(\ref{fig:my_label1})(b). In the Fig.~(\ref{fig:my_label4}), we can also find the same behavior for the entropy densities in both 3D and 2D space though there is a slight increase in the ratio for non-relativistic case compared to other cases.

%
    \begin{figure}
    \centering
    \includegraphics[scale=0.25]{Fig2_Gen_eta_etasb_muT_3D.eps}
    \includegraphics[scale=0.25]{Fig2_Gen_eta_etasb_muT_2D.eps}
    \caption{The ratio of shear viscosity of electron flow in different domains to $\eta_{SB}$ with $\mu/T$ (a) in the 3D case and (b) in the 2D case respectively}
    \label{fig:e}
    \end{figure}
%
    Now, let us come to the shear viscosity to entropy density ratio $\eta/s$, which is a more important quantity than only $\eta$ to measure the fluidity of the system. In the D.F domain, the extreme situation (mathematically) is $\mu \rightarrow 0$. In this limit, $\eta$ and $s$ carry quite similar terms, so when we take their ratio, we will get very simplified expressions:
    \bea 
    \frac{\eta}{s} &=& \frac{\tau_c T}{5} ~{\rm for ~3D ~URF/GF~,}\\
    &=& \frac{\tau_c T}{10} ~{\rm for ~3D ~NRF~,}\\
    &=& \frac{\tau_c T}{4} ~{\rm for ~2D ~URF/GF~,}\\
    &=& \frac{\tau_c T}{8} ~{\rm for ~2D ~NRF~.}
    \eea
    From the String theory-based calculation~\cite{25Kovtun:2004de}, it was conjectured that $\eta/s$ has a lower bound, well known as KSS bound, which gives an inequality $\frac{\eta}{s} \geq \frac{\hbar}{k_B} \frac{1}{4 \pi}=\frac{1}{4 \pi} $ (in natural unit). Though classically one may expect $\tau_c \rightarrow 0 \Rightarrow \frac{\eta}{s} \rightarrow 0$, but quantum mechanically, relaxation or scattering time $\tau_c$ or mean free path $\lambda_c \approx v \tau_c$ can not be lower than de-Brogile range of time or wavelength scale. This simple quantum mechanical concept also suggests a lower bound of $\frac{\eta}{s}$, sometimes called a quantum lower bound. By imposing this bound $\frac{\eta}{s}= \frac{1}{4 \pi}$, we can get a rough expression lower bound of $\tau_c$ as,
    \bea 
    \tau_c &=& \frac{5}{4 \pi T} ~{\rm for ~3D ~URF/GF~,}\\
    &=& \frac{10}{4 \pi T} ~{\rm for ~3D ~NRF~,}\\
    &=& \frac{4}{4 \pi T} ~{\rm for ~2D ~URF/GF~,}\\
    &=& \frac{8}{4 \pi T} ~{\rm for ~2D ~NRF~.}
    \label{tc_KSS}
    \eea
%
    \begin{figure} %[ht]
    \centering
    \includegraphics[scale=0.25]{Fig6A_Gen_etas_T_2D.eps}
    \includegraphics[scale=0.25]{Fig8B_Gen_etas_muT_2D.eps}
    \caption{Shear viscosity to entropy density ratio; (a) with $T$ for $\mu=0$ (Undoped graphene) and (b) with $\mu/T$}
    \label{fig:es}
    \end{figure}
%
    This KSS bound conjecture~\cite{25Kovtun:2004de} makes the scientific community curious to find such fluid whose $\eta/s$ is close to this bound. In other words, if we write $\eta/s=n/(4\pi)$, where $n\geq 1$~ \footnote{$n$ should not be confused with number density}, then fluid with $n=1-5$ may be considered as those special fluids and may be called nearly or close to perfect fluid. Empirically, QGP is the evidence of such perfect fluid ($n\approx 1-2$) in the relativistic domain, while close to perfect fluid ($n\approx 5$) example for NR case is cold atom systems~\cite{27Sch_fer_2009}. According to Eqs.~(\ref{tc_KSS}), we can expect gross values of relaxation time for QGP and cold atom systems as $\tau_c\approx\frac{5}{4\pi T}$-$\frac{10}{4\pi T}$ and $\tau_c\approx\frac{50}{4\pi T}$ respectively. 
    Similarly, according to the theoretical prediction from Ref.~\cite{24PhysRevLett.103.025301}, GF may also belong to this close-to-perfect fluid category. So far, to the best of our knowledge, no experimental measurement of $\eta/s$ vs. $T$ plot is available, so the theoretical plot of $\eta/s$ vs. $T$ in Ref.~\cite{24PhysRevLett.103.025301} is considered as our reference to guess or tune order of magnitude for $\tau_c\approx\frac{n}{\pi T}$. In the left panel of Fig.~(\ref{fig:es}), we can get guidance that $\tau_c\approx\frac{n}{\pi T}$ within $n=3$-$5$ can cover the order of magnitude of $\eta/s$ in the temperature range $T=35$-$150^\circ$K, predicted by Müller et al.~~\cite{24PhysRevLett.103.025301}. By considering an average value $\tau_c\approx\frac{4}{\pi T}$, we have plotted $\eta/s$ of 2D GF or URF (red solid line) and 2D NRF (green dash line) against $\mu/T$-axis in the right panel of Fig.~(\ref{fig:es}).
    We notice that $\eta/s$ in DF domain becomes lower than in the FL domain, mainly because of the thermodynamical phase-space part of $\eta/s$. In terms of Fermi integral function, this part for GF can be identified as $\Big[3f_3\left(A\right) - \frac{\mu}{T} f_2\left(A\right) \Big]^{-1} f_3\left(A\right)$ from Eq.~(\ref{2dg}).
    We have put the NRF case for reference, but 2D NRF for electrons may be possible in a hypothetical situation. So present study indicates that dropping $\eta/s$ values and saturating towards constant values may be found during the transition from FL to DF domains in the graphene system.
    Although we have a limitation in that we have considered $\tau_c\propto 1/T$, which may be changed in the actual microscopic calculation of $\tau_c$, and so the trend of $\eta/s$ may also be changed. It demands more theoretical studies on these $\eta/s$ estimations as well as the explicit measurement of this quantity from the experimental side.


    %After knowing the lower bound conjecture of shear viscosity to entropy density $(\eta/s)$ as $\hbar/(4\pi k_B)$ or $1/(4\pi)$ (in natural unit)~\cite{25Kovtun:2004de}, scientific communities are curious to know those strongly coupled systems, which are close to that bounds. Experimentally, the RF, like quark and hadronic matter, produced in high energy heavy ion collision experiments and NRF like cold atom systems~\cite{27Sch_fer_2009} are identified as those strongly coupled systems. On the other hand, GF may also belong to that category according to the theoretical prediction from Ref.~\cite{24PhysRevLett.103.025301}, which is considered as reference point for tuning our results. 

%
    %The lower KSS value $1/4\pi$ from the experimental result of RHIC and LHC matter \cite{27Sch_fer_2009} is less than the theoretical value of $\eta/s=\tau_c T/5$ in quark matter. Consequently, the corresponding relaxation time for QGP matter will be $\tau_c=5/4\pi T$. However, the theoretical expression for the graphene system is $\eta/s=\tau_c T/4$. Accordingly, the relaxation time can be taken as $\tau_c=1/\pi T$ obtained in correlation with the lower bound of QGP matter. In the Fig.~(\ref{fig:my_label1})(a), the maroon line \cite{24PhysRevLett.103.025301} is shown as a reference from the calculation using the quantum kinetic theory for clean undoped graphene. In our case, $\tau_c$ is considered a free parameter to facilitate the comparison of the result in Ref.~\cite{24PhysRevLett.103.025301}. It means that it can be any arbitrary value. Therefore, the general expression of $\tau_c$ can be regarded as $\tau_c=n/\pi T$. We have taken the relaxation time of electrons in the graphene system as $\tau_c=3/\pi T$ and $\tau_c=5/\pi T$, respectively, three times and five times larger than the relaxation time of the corresponding lower bound of quark matter. These values are saturated. Our shear viscosity to entropy density ratio with $n=5$ from the relaxation time approximation method satisfies the Muller data \cite{24PhysRevLett.103.025301} at $T=35 K$. This ratio with $n=3$ satisfies $T=150 K$, respectively. Comparing these two $\tau_c$ values using the Eq.~\ref{2dg}, we can say that the smaller $\tau_c$ value corresponds to the smaller $\eta/s$ such that there are more chances to find the Hydrodynamics flow in electron fluid on Graphene. So, our approach, the relaxation time approximation based Boltzmann equation, to get the shear viscosity calculation is one of the appropriate ways for the appropriate $\eta/s$ close to the KSS bound.
%
    %By tuning the value of $\tau_c=4/\pi T$, which is four times larger than the KSS bound, the value of shear viscosity to entropy density ratio has been evaluated and plotted with $\mu/T$ at specific temperature $T=60 K$ in Fig.~(\ref{fig:my_label1})(b), represented by the solid red line for graphene and solid green line for non-relativistic case. The corresponding Sommerfeld-Lemma values of the two cases are also included to see the variation in the actual value of the ratio.
    %It is said that there are two regions; one below $\mu/T\approx 1$ is the Dirac Fluid region(D.F), and another beyond $\mu/T\approx 1$ is the Fermi Liquid region(F.L) in Ref.~\cite{?}. 
    %The D.F is also known as a fluid in which electron-electron scattering is very fast and where electron hydrodynamics dominate; other electron-phonon and electron-carrier interactions are significantly less Ref.~\cite{17doi:10.1126/science.aad0343}. Realization of the Dirac fluid in graphene requires that the thermal energy be larger than the local chemical potential $\mu(r)$, defined at position $r: k_B T \geq |\mu(r)|$ \cite{17doi:10.1126/science.aad0343}. In the relaxation time scenario, fewer carriers or impurities in the D.F region means fewer interactions occur among these constituents. It also means that these fewer impurities do not affect the fluid flow of electrons in Graphene and non-relativistic cases. As a result, the relaxation time will be less in the D.F region of Graphene since it is directly proportional to the mean free path of the particles $\lambda= v.\tau_c$ ($v$ is the electron velocity). Therefore,  $\eta/s$ is close to the KSS bound of quark matter. In contrast, the ratio increases in the F.L region; it shows more interactions between the electrons, phonons, and carriers because some impurities and phonons can occur in this F.L. In practice, the relaxation time for this region will be different due to more interaction processes. Still, we used the same $\tau_c$ in calculating $\eta/s$ in this F.L region as in the D.F region. 

    %On the other hand, there is no patent change in the value of the ratio below $\mu/T\approx 1$, but there is a slight change in shear viscosity to entropy ratio at the temperature of $T=60\ K$ with some low $\mu$ values. The same behavior of the electron flow is also found in the non-relativistic case within the low chemical potential to temperature value in the D.F region. On the other hand, the ratio values deviate away from the constant values and merge as a straight line in the Fermi Liquid region(F.L). It can be said that the same fluid property of electrons on graphene can be seen in the non-relativistic case.  
%
 %\begin{figure}
 %    \centering
 %    \includegraphics[scale=0.25]{Fig1_Gen_eta_muT_3D.eps}
 %    \includegraphics[scale=0.25]{Fig1_Gen_eta_muT_2D.eps}
 %    \caption{A. Normalized shear viscosity of electron flow different domains with $\mu/T$ in the 3D case and, B. Normalised shear viscosity of different domains with $\mu/T$ in 2D case}
 %    \label{fig:my_label2}
 %\end{figure}
%
%
    %Some of the interesting explorations regarding graphene are the normalized shear10 viscosity in the low-temperature case for chemical potential zero and the normalized shear viscosity at absolute zero temperature case. The normalized values of shear viscosity for $\mu=0$ in graphene and ultra-relativistic cases are constant in 3D space. That value in 2D space is decreased with temperature in Fig.~(\ref{fig:my_label5}). We noticed that the decrease in $\eta$s is not obvious in 2D space according to the order of normalized shear viscosity. The value for the non-relativistic case is decreased with temperature in both 3D and 2D space. For absolute zero temperature, we can make the same discussion for the increase in chemical potential in Fig.~(\ref{fig:my_label6}).
%
    \section{Summary and Conclusion}
    \label{sec: Sum}
%    Throughout this article, we have addressed the hydrodynamic description of the Graphene system in terms of the energy-momentum tensor. 
    We can summarize our investigation in the following steps.
    First, we introduce a brief macroscopic description of electron fluid in graphene, then we focus on its microscopic description. Our dealing quantity is considered as energy-momentum tensor, whose ideal part represents energy density and pressure in the static limit picture of fluid dynamics. Using those thermodynamical quantities, our destination from the ideal part of the energy-momentum tensor becomes entropy density, which will be used to be normalized with shear viscosity. From the dissipative part of the energy-momentum tensor, shear viscosity coefficients of electron fluid are calculated based on the kinetic theory approach with relaxation time approximation. Temperature and chemical potential-dependent general expressions of shear viscosity, entropy density, and the shear viscosity to entropy density ratio has been calculated and plotted for different cases of electron fluid like non-relativistic, graphene and ultra-relativistic cases. For completeness of comparison, we considered both 3D and 2D systems. Analyzing the results of different cases, we get a comparative understanding and conclusions, which are addressed briefly in bullet points:
    \begin{itemize}
    \item $\mu/T$ dependence of shear viscosity $\eta$ as well as entropy density $s$ for URF and GF are exactly similar but a little different from the NRF.
    \item We notice a huge difference among URF, GF, and NRF in terms of the order of magnitude of $\eta$ and $s$ with ranking URF $\ll$ GF $\ll$ NRF.
    \item During transiting from 3D to 2D, order of magnitude of $\eta$ and $s$ shift towards lower values
    \item When we go from Fermi Liquid ($\mu/T\gg 1$) to Dirac Liquid ($\mu/T\ll 1$) domain, values of $\eta$,
        $s$ and $\eta/s$ ratio decrease towards a saturated values.
    \item Interesting ranking for $\eta/s$ become  URF $=$ GF $\ge$ NRF
    \end{itemize}
    Present comparative study on the microscopic calculation of shear viscosity may be considered as a good documentation of master formulas of different cases from 3D URF, GF, NRF to 2D URF, GF and NRF. In the future, it may be useful for actual graphene system estimation, where one should go with some first principle or model dependent calculation of relaxation time. Also, one should deal with electron-hole plasma with appropriate degeneracy factor in the Dirac Fluid domain for the actual graphene system but the present work sticks with electron description only due to observing the estimations for its  different dispersion relations. Our immediate future plan is to concentrate on that actual graphene phenomenology on the viscous aspects. 
    
    %Taking the $n=3-5$ in the expression of $\tau_c=n/\pi T$, and comparing the resultant $\eta/s$, we explained the dependence of $\tau_c$ on the hydrodynamic flow of electrons on graphene and our result of $\eta/s$ satisfies the existing theoretical result. By taking the value $\tau_c=4/\pi T$, our calculated value $\eta/s$ can explain why the existence of the electron hydrodynamics in D.F affects the interactions of the number of impurities in the larger thermal energy than the chemical potential. Furthermore, the ratio differs in the F.L region, where the thermal energy is less than the chemical potential.

    %It is found that the ratio band covers the theoretical value from the Ref.~\cite{24PhysRevLett.103.025301} and closes to the KSS bound. For $n=4$, the ratio said that the flow of electrons in graphene has more fluidity at low $\mu/T$ in the Dirac fluid (D.F) and less fluidity in the Fermi liquid (F.L) domain. The normalized shear viscosity also increases beyond $\mu/T\approx 1$ with their different orders of shear viscosity values. The ratio of the shear viscosity $\eta/\eta_{SB}$ has also increased beyond the boundary $\mu/T\approx 1$ in all cases. The same behavior of fluid property has happened for the entropy densities.
    %\textcolor{red}{Pending:: We will try to explain why every fluid property of all the cases has some changes in the boundary beyond $\mu/T\approx 1$.}
%%%%%%%%%%%%%%%%%%%%%%%%%%%%%%%%%%%%%%%%%%%%%%%%%%%%%%%%%%
%
%    \includegraphics[scale=0.25]{}
%    \includegraphics[scale=0.25]{Fig9_Gen_eta_muT_2D.eps}
%    \caption{3D General Case later and 2D General case}
%    \includegraphics[scale=0.25]{}
%    \includegraphics[scale=0.25]{Fig10_Gen_s_muT_2D.eps}
%    \caption{3D General Case later and 2D General case}
%\end{figure}
%
%\begin{figure}
%    \centering
%    \includegraphics[scale=0.25]{Fig1_Gen_eta_muT_3D.eps}
%    \includegraphics[scale=0.25]{Fig1_Gen_eta_muT_2D.eps}
%    \caption{2D General Case}
%    \label{fig:my_label}
%\end{figure}
%
%\begin{figure}
 %   \centering
  %  \includegraphics[scale=0.25]{Th_s_T_m.eps}
   % \caption{Entropy density with temperature for the case of $\mu= 0$}
%    \label{fig:my_label}
%    \includegraphics[scale=0.25]{eta_mu_a25nm.eps}
%    \includegraphics[scale=0.25]{eta_T_a25nm.eps}
%    \caption{Normalised shear viscosity of different descriptions with chemical potential and temperature respectively: (a) the variation of shear viscosity with chemical potential for the specific case of T=0 and (b) the variation with temperature for the case of $\mu= 0$}
%    \label{fig:my_label}
%    \includegraphics[scale=0.25]{tc_mu_a25nm.eps}
%    \includegraphics[scale=0.25]{tc_T_a25nm.eps}
%    \caption{Relaxation time of different descriptions with chemical potential and temperature respectively: (a) the variation of relaxation time with the chemical potential for the specific case of T=0 and (b) the variation with temperature for the case of $\mu= 0$}
%    \label{fig:my_label}
%   \includegraphics[scale=0.25]{eta_by_s_T_a25nm.eps}
%    \caption{Shear viscosity to entropy ratio of different descriptions with temperature compared with the KSS bound of QGP medium}
%    \label{fig:my_label}
%\end{figure}
%\newpage
    \begin{acknowledgments}
    This work was partly (CWA and TZW) supported by the Doctoral Fellowship in India (DIA) program of the Ministry of Education, Government of India. The authors thank the other members of eHD club - Sesha P. Vempati, Ashutosh Dwibedi, Narayan Prasad, Bharat Kukkar, and Subhalaxmi Das Nayak.
    \end{acknowledgments} 
%
%\hfill \break
    \appendix
    \section{Density of States}
    \label{App_A}
%\hfill \break
%\begin{center}
%    \textbf{\uppercase{B.}}
%    \textbf{\uppercase{\underline{Density of States}}}
%\end{center}
    The density of states is nothing but the total number of energy states per unit energy interval. If the total number of energy states in energy range $E$ to $E + dE$ are $D\left(E\right)dE$, then the density of states will be 
    \begin{equation}
    g\left(E\right) = \frac{D\left(E\right)dE}{dE}~.
    \end{equation}
    Now, here are some expressions of the number of energy states for different-different cases expressed as

    \hfill \break
    \textbf{\underline{Case:1.}} for 3D Graphene,
    \begin{equation}
    D\left(E\right)dE = N_s \frac{4\pi V}{h^3 v_g^3}E^2 dE~,
    \end{equation}
    \textbf{\underline{Case:2.}} for 3D Non-Relativistic,
    \begin{equation}
     D\left(E\right)dE = N_s 2\pi V \left(\frac{2m}{h^2}\right)^{\frac{3}{2}}\sqrt{E} dE~,
    \end{equation}
    \textbf{\underline{Case:3.}} for 2D Graphene,
    \begin{equation}
    D\left(E\right)dE = N_s \frac{2\pi S}{h^2 v_g^2} E \, dE~,
    \end{equation}
    \textbf{\underline{Case:4.}} for 2D Non-Relativistic,
    \begin{equation}
     D\left(E\right)dE = N_s \frac{2\pi S}{h^2} m\, dE~.
    \end{equation}
    In the above expressions, $V$ and $S$ represent the volume and area in position space, respectively.

% \hfill \break
    \section{Fermi-Dirac Integral}
    \label{App_B}
%\hfill \break
%\begin{center}
%    \textbf{\uppercase{\underline{C. Fermi-Dirac Integral}}}
%\end{center}
    We have the integral form as
    \begin{align}
    f_\nu (A)=\frac{1}{\Gamma (\nu)}\int_0^\infty \frac{x^{\nu-1}}{A^{-1} e^x+1} dx ~,
    \label{mom}
    \end{align}
    where $f_\nu (A)$ is known as the Fermi-Dirac integral and $x= \beta E$. And the expression in terms of energy can be written in terms of $x$ as
    \begin{align}
    \int_{0}^{\infty}  \frac{E^{\nu -1}}{A^{-1}e^{\beta E} + 1} dE = & \frac{1}{\beta^\nu} \int_0^\infty \frac{x^{\nu-1}}{A^{-1} e^x+1} dx = \frac{1}{\beta^\nu} \Gamma (\nu)  f_\nu (A)~.
    \end{align}
%
% \hfill \break
% \section{Sommerfeld-Lemma Expressions 2D-Case Graphene}
% \label{App_C}
%\hfill \break
%\hfill \break
%\begin{center}
%    \textbf{\uppercase{D.}}
%    \textbf{\uppercase{\underline{Sommerfeld-Lemma Expressions 2D-Case Graphene}}}
%\end{center}
% \textbf{\underline{Case:1. 2D Graphene}}
% \begin{align}
%    s_g^{2D} = \frac{\pi}{3 v_g^2} \mu T
% \end{align}
% \begin{align}
%    \eta_g^{2D} = \frac{N_s}{16 \pi v_g^2} \tau_c \mu^3 \left[ 1+\pi^2 \left( \frac{T}{\mu} \right)^2 \right]
% \end{align}
% \begin{align}
%   \frac{\eta_g^{2D}}{s_g^{2D}} = \frac{3 N_s }{16 \pi^2} \tau_c \left[ 1 + \pi^2 \left(\frac{T}{\mu}\right)^2 \right] \frac{\mu^2}{T}
%   \label{mom}
% \end{align}
%
% \textbf{\underline{Case:2. 2D Non-relativistic}}
% \begin{align}
%    s_{NR}^{2D} = \frac{m \pi}{3} T
% \end{align}
% \begin{align}
%    \eta_{NR}^{2D} = \frac{N_s m }{16 \pi} \tau_c \mu^2 \left[ 1+\frac{\pi^2}{3} \left( \frac{T}{\mu} \right)^2 \right]
% \end{align}
% \begin{align}
%   \frac{\eta_{NR}^{2D}}{s_{NR}^{2D}} = \frac{3 N_s }{16 \pi^2} \tau_c \left[ 1 + \frac{\pi^2}{3} \left(\frac{T}{\mu}\right)^2 \right] \frac{\mu^2}{T}
%   \label{mom}
% \end{align}
%
% \hfill \break
    \section{Average Angular Integral in 2D}
    \label{App_D}
%\begin{center}
%     \textbf{\uppercase{\underline{E. Average Angular Integral in 2D}}}
%\end{center}
    We have the integral form as
    \begin{align*}
    \int p_i p_j p_k p_l \, d^2p = p dp \int  p_i p_j p_k p_l \,d\theta~.
    \end{align*}
    Since
    \begin{align*}
    \Vec{p} = p \hat{n}
    \end{align*}
    where, $$ \hat{n} = \cos{\theta}\,\hat{i} + \sin{\theta}\, \hat{j}~, $$
    and
    $$ p_i = \Vec{p}.\hat{e_i} = p \left(\hat{n}.\hat{e_i}\right) = pn_i~, $$
    the integral becomes
    \begin{align*}
     \int p_i p_j p_k p_l \, d^2p = p\,dp \, p^4 \int  n_i n_j n_k n_l \,d\theta~.
    \end{align*}
    Now, we have to calculate
    $$ \int  n_i n_j n_k n_l \,d\theta = ? $$
    \underline{\textbf{Case:1}} %When $i = j = 1$ and $k = l = 2$\\
    The above integral becomes
    \begin{align*}
    \int n_1^2 n_2^2\, d\theta &= \int_0^{2 \pi} \cos^2{\theta} \sin^2{\theta} \, d\theta~,\\
    & = 4  \int_0^{ \frac{\pi}{2}} \cos^2{\theta} \sin^2{\theta} \, d\theta~.
    \end{align*}
    Now, using the Beta function, we know that
    \begin{align}
   & B\left( u, v\right) = 2 \int_0^{\frac{\pi}{2}} \left(\cos{\theta}\right)^{2u-1}\left( \sin{\theta}\right)^{2v-1} \, d\theta~,\\
   \implies & B\left( u, v\right) = \frac{\Gamma u\, \Gamma v}{\Gamma\left(u+v\right)}~.
    \end{align}
    Applying this, we get
    \begin{equation}
    \int n_1^2 n_2^2\, d\theta = \frac{2 \pi}{8}~.
    \end{equation}
    \underline{\textbf{Case:2}}% When $i = j = k = 1$ and $l = 2$  \\
    \begin{align*}
    \int n_1^3 n_2\, d\theta &= \int n_1 n_2^3\, d\theta~,\\
    & = \int_0^{2 \pi} \cos^3{\theta} \sin{\theta} \, d\theta = 0~.
    \end{align*}
    \underline{\textbf{Case:3}}\\
    \begin{align*}
    \int n_1^4\, d\theta &=  \int n_2^4\, d\theta~, \\
    & =  \int_0^{2 \pi} \sin^4{\theta} \, d\theta~,\\ 
    &= 2 \, B\left( \frac{5}{2}, \frac{1}{2}\right) = \frac{3 \pi}{4}~.
    \end{align*}
    Now, the above integral can be written as
    \begin{align*}
     \int n_1^4\, d\theta &=  \int n_2^4\, d\theta = \frac{2 \pi}{8}\times 3~,
    \end{align*}
    Now, the integral can be expressed as
    \begin{align}
    \int  n_i n_j n_k n_l \,d\theta = \frac{2\pi }{8} \left( \delta_{ij}\delta_{kl} + \delta_{ik}\delta_{jl} + \delta_{il}\delta_{jk}\right)~,
    \end{align}
    and
    \begin{align*}
     \int p_i p_j p_k p_l \, d^2p = 2\pi \, p\,dp \frac{p^4}{8}\left( \delta_{ij}\delta_{kl} + \delta_{ik}\delta_{jl} + \delta_{il}\delta_{jk}\right)~.
    \end{align*}
    Now, the average final expression will be
    \begin{equation}
    <p_i p_j p_k p_l> = \frac{p^4}{8}\left( \delta_{ij}\delta_{kl} + \delta_{ik}\delta_{jl} + \delta_{il}\delta_{jk}\right)~.
    \end{equation}
%\include{Appendix}
%\newpage
%\begin{thebibliography}{999}
%	\input{bib.tex}
%\end{thebibliography}
    \bibliography{new_ref}
%\bibliography{}
    \end{document}
